\ifdefined\MyWebBook\else
\documentclass[../textbook.tex]{subfiles}
\begin{document}
\fi
\bookchapter{人类反馈学习}{ch:preference-alignment}

监督微调要求示范者给出应当生成的回答，偏好数据则允许示范者在多个候选之间作比较。这种差异使反馈更容易针对语气、完整性或任务适配程度，但也引入了新的统计问题：一次比较究竟代表谁的偏好，候选从哪里产生，以及局部排序能否外推到训练后出现的新回答。

比较反馈可表示为概率模型；将其与奖励正则化目标结合，可以得到直接偏好优化。在线强化学习的策略梯度基础见\bookxref{ch:reinforcement-learning}；可验证奖励和推理训练在\bookxref{ch:verifiable-reasoning}进一步展开。这里的“对齐”指依据给定反馈分布调整模型行为。

\section{人类反馈与偏好数据}

\subsection{比较单位与数据分布}
一条偏好记录可表示为 $(x,y_a,y_b,z)$，其中 $x$ 为共同上下文，$y_a,y_b$ 为两个候选，$z\in\{0,1\}$ 指示是否偏好 $y_a$。上下文包括角色、约束和必要历史；若两条候选实际上回答了不同问题，同一提示下的偏好模型便解释不了它们。

\term{偏好对}{Preference Pair}{}的生成机制至少包含提示分布、候选策略、解码配置和标注规则。设提示来自 $\mu$，候选来自条件分布 $q(y_a,y_b\mid x)$，标注者来自群体分布 $h$，则训练风险的期望同时依赖这三者。只扩大比较数量，不扩大提示和候选覆盖，可能只是反复学习同一狭窄区域。

候选显示次序、长度、格式、模型身份泄漏和标注者知识都会影响比较。应保存可重建的原始候选、匿名化呈现条件和标注版本。多个候选产生的成对组合共享同一提示和回答，彼此相关；划分训练与评估集合时应以提示或更高层来源分组。

\subsection{符号与假设}
本章主要符号见表\ref{tab:preference-alignment-symbols}。
\begin{table}[htbp]
\centering
\caption{偏好学习的主要符号。}
\label{tab:preference-alignment-symbols}
\begin{tabularx}{\linewidth}{@{}lX@{}}
\toprule
符号 & 含义\\
\midrule
$x,y,T_y$ & 提示、回答序列和回答所计词元数，终止标记是否计入须固定。\\
$y_w,y_l$ & 标注中相对优选与相对拒选的回答。\\
$r_{\vect{\phi}}(x,y)$ & 参数为 $\vect{\phi}$ 的标量奖励函数。\\
$\pi_{\vect{\theta}},\pi_{\mathrm{ref}},\pi_{\mathrm{old}}$ & 当前策略、固定参考策略及产生一批在线样本的旧策略。\\
$\beta$ & KL 正则化系数；DPO 推导沿用同一尺度。\\
$\sigma(u)$ & Logistic 函数 $\frac{1}{1+\conste{}^{-u}}$。\\
$Z(x)$ & 固定提示下的归一化配分函数。\\
$\Delta_{\vect{\theta}}$ & 优选与拒选回答的相对参考对数概率差。\\
\bottomrule
\end{tabularx}
\end{table}

\begin{assumption}[偏好与策略推导的范围]
先假设固定提示下回答集合有限，参考策略在该集合上为正，奖励有限且 $\beta>0$。偏好用一个标量效用差的 Logistic 概率描述。对可数无限回答集合，还需保证配分函数有限及相关期望存在。真实群体可能有冲突、平局和上下文依赖，单标量模型只是其中一种近似。
\end{assumption}

\section{奖励建模}

\subsection{Bradley--Terry 形式}
将每个候选的正强度写为 $u_a=\conste{}^{r_a}$ 和 $u_b=\conste{}^{r_b}$，以相对强度定义胜出概率，得到 Bradley--Terry 形式\autocite{bradley1952paired}
\begin{equation}
 p_{\vect{\phi}}(y_a\succ y_b\mid x)
 =\frac{\conste{}^{r_a}}{\conste{}^{r_a}+\conste{}^{r_b}}
 =\sigma(r_a-r_b),\qquad r_a=r_{\vect{\phi}}(x,y_a).
 \label{eq:pref-bt}
\end{equation}
对二元标签采用负对数似然，单对损失为
\begin{equation}
 \ell_{\mathrm{RM}}=-z\log\sigma(d)-(1-z)\log\sigma(-d),
 \qquad d=r_a-r_b.
 \label{eq:pref-rm-loss}
\end{equation}
其间隔导数为 $\frac{\partial\ell}{\partial d}=\sigma(d)-z$；二阶导数为 $\sigma(d)(1-\sigma(d))\ge0$。损失对间隔凸，深度奖励网络的参数优化通常仍为非凸问题。标注偏好 $a$ 却预测 $\sigma(d)$ 很小时，梯度会推动 $r_a-r_b$ 增大。

如果同一对有多个独立判断，$z$ 可替换为偏好 $a$ 的比例，表示二项似然按比较数归一化后的目标。把一个平局直接写为 $z=\frac{1}{2}$，相当于要求模型对两者给出相等偏好概率；它与带显式平局概率的三分类机制含义不同，应根据标注含义选择。

\subsection{奖励的可辨识性}
对同一提示给所有奖励加上任意 $c(x)$，式\eqref{eq:pref-bt}不变。因此成对比较无法识别提示相关的绝对零点。若数据只连接一部分回答子集，不同连通分量之间的相对偏移也无法由这些比较确定。

尺度则不同：将差值乘常数会改变偏好概率。若同时引入未知噪声温度 $\tau$，模型只由 $\frac{r_a-r_b}{\tau}$ 决定，此时奖励尺度和温度混在一起，无法分别识别。奖励归一化或温度修改会改变后续策略优化的有效强度，应与策略目标一并记录。

标量效用还不能表示确定性的循环偏好：若 $a\succ b,b\succ c,c\succ a$ 均被要求以大于 $\frac{1}{2}$ 的概率成立，则分别需要 $r_a>r_b>r_c>r_a$，产生矛盾。出现系统性冲突时，应检查群体、任务和呈现条件；增加训练轮数解决不了这种矛盾。

\subsection{奖励网络与校准}
奖励模型可以用编码器或因果模型提取 $(x,y)$ 的表示，再经标量读出头输出 $r_{\vect{\phi}}$。若使用回答末端表示，须取实际有效末端而不是填充位置；若截断删除了结论或约束，原来的偏好标签可能不再适用。模型结构并不要求与策略完全相同，但输入语义必须兼容。

排序准确率关注差值符号，概率校准关注预测胜率与观察频率是否一致，两者测的是不同方面。即使在离线比较集上校准良好，策略优化后产生的新回答也可能偏离原候选分布。对奖励模型进行优化会主动寻找其高分区域，因此评估需要覆盖被优化后的候选，而不只是原始数据。

设$r_w=1.2,r_l=0.2$，间隔为1，预测胜率 $\sigma(1)\approx0.7311$，优选标签损失约0.3133。若在相同条件下只有55\%的人选择 $w$，预测胜率比观察比例高约18个百分点。换成均为0.7的奖励，预测胜率为0.5，表示两者具有相同效用。

\section{基于人类反馈的强化学习}

\subsection{目标与 KL 方向}
\term{基于人类反馈的强化学习}{Reinforcement Learning from Human Feedback}{RLHF}可通过比较数据拟合奖励，再优化策略；InstructGPT 是这种训练链路的代表性工作\autocite{ouyang2022instructgpt}。固定提示下的一种目标为
\begin{equation}
 J_x(\pi)=\sum_y\pi(y\mid x)r(x,y)
 -\beta\sum_y\pi(y\mid x)\log\frac{\pi(y\mid x)}{\pi_{\mathrm{ref}}(y\mid x)}.
 \label{eq:pref-regularized}
\end{equation}
KL 的方向为当前策略到参考策略。它使参考零概率事件具有无穷偏离代价，所以支持集条件是推导的一部分。参考约束控制分布漂移，奖励可靠性、事实准确性和原有能力通过任务评估检查。

\subsection{最优策略的推导}
对归一化约束 $\sum_y\pi_y=1$ 引入乘子 $\lambda$，记 $q_y=\pi_{\mathrm{ref}}(y\mid x)$，拉格朗日函数为
\begin{equation}
 \mathcal F=\sum_y\pi_yr_y-\beta\sum_y\pi_y\log(\frac{\pi_y}{q_y})
 +\lambda(\sum_y\pi_y-1).
\end{equation}
对正概率内部点求导得
\begin{equation}
 r_y-\beta[\log(\frac{\pi_y}{q_y})+1]+\lambda=0.
\end{equation}
整理可得 $\pi_y=q_y\exp(\frac{r_y}{\beta})\exp(\frac{\lambda}{\beta}-1)$。最后一项不依赖回答，通过归一化确定为 $\frac{1}{Z(x)}$，因此
\begin{equation}
 \pi^*(y\mid x)=\frac{\pi_{\mathrm{ref}}(y\mid x)\conste{}^{\frac{r(x,y)}{\beta}}}{Z(x)},
 \qquad Z(x)=\sum_y\pi_{\mathrm{ref}}(y\mid x)\conste{}^{\frac{r(x,y)}{\beta}}.
 \label{eq:pref-optimal-policy}
\end{equation}
不仅可以检查驻点，还可以直接证明全局最优。把最优策略的对数表达代回式\eqref{eq:pref-regularized}，得到
\begin{equation}
 J_x(\pi)=\beta\log Z(x)-\beta D_{\mathrm{KL}}(\pi\|\pi^*).
 \label{eq:pref-optimal-proof}
\end{equation}
KL 非负，且仅在两个分布相等时为零，所以 $\pi^*$ 是支持集内的唯一最优分布。这个证明同时给出任意策略距离最优目标值的差额。

当固定奖励且 $\beta$ 增大时，指数倾斜减弱；当 $\beta\to\infty$ 时趋近参考策略。$\beta\to0^+$ 时质量集中到最高奖励集合，并在并列最优集合中保留参考权重比例。奖励与 $\beta$ 同时乘同一个正数不改变式\eqref{eq:pref-optimal-policy}，再次说明奖励尺度与约束强度必须一起解释。

\begin{example}


设回答集合仅有 $a,b$，参考概率为 $(0.8,0.2)$，奖励为 $(0,\log4)$，$\beta=1$。未归一化质量为 $(0.8,0.8)$，$Z=1.6$，故最优策略为 $(0.5,0.5)$。

最优期望奖励为 $\frac12\log4=\log2\approx0.6931$。KL 为
\begin{equation*}
 \frac12\log\frac{0.5}{0.8}+\frac12\log\frac{0.5}{0.2}
 =\frac12\log1.5625\approx0.2231.
\end{equation*}
目标值为0.4700，等于 $\log1.6$。虽然 $b$ 的奖励较高，最优策略并未以概率1生成它；参考约束保留了概率分配的代价。若将两奖励同时增加10，最优策略不变，而目标值增加10。
\end{example}

\section{直接偏好优化}

\subsection{由最优策略表示奖励}
式\eqref{eq:pref-optimal-policy}等价于
\begin{equation}
 r(x,y)=\beta\log\frac{\pi^*(y\mid x)}{\pi_{\mathrm{ref}}(y\mid x)}+\beta\log Z(x).
\end{equation}
同一提示下，两回答奖励相减时，配分函数项消失。将策略参数化为 $\pi_{\vect{\theta}}$，用相对参考的对数概率差表示偏好间隔，就得到\term{直接偏好优化}{Direct Preference Optimization}{DPO}\autocite{rafailov2023dpo}：
\begin{align}
 \Delta_{\vect{\theta}}&=\log\pi_{\vect{\theta}}(y_w\mid x)-\log\pi_{\vect{\theta}}(y_l\mid x)
 -\log\pi_{\mathrm{ref}}(y_w\mid x)+\log\pi_{\mathrm{ref}}(y_l\mid x),\\
 \ell_{\mathrm{DPO}}&=-\log\sigma(\beta\Delta_{\vect{\theta}})
 =\operatorname{softplus}(-\beta\Delta_{\vect{\theta}}).
 \label{eq:pref-dpo}
\end{align}
这将显式奖励模型的拟合转成策略上的偏好似然。实际策略与式\eqref{eq:pref-optimal-policy}的理想解之间仍有数据覆盖、模型容量和优化误差。

\subsection{偏好梯度与概率变化}
设 $u=\beta\Delta_{\vect{\theta}}$，则
\begin{equation}
 \nabla_{\vect{\theta}}\ell_{\mathrm{DPO}}
 =-\beta\sigma(-u)
 [\nabla_{\vect{\theta}}\log\pi_{\vect{\theta}}(y_w\mid x)
 -\nabla_{\vect{\theta}}\log\pi_{\vect{\theta}}(y_l\mid x)].
 \label{eq:pref-dpo-gradient}
\end{equation}
参考概率固定，不参与求导。排序错误且间隔很负时，系数趋近 $-\beta$；间隔很正时，梯度趋近零。$\beta$ 同时改变间隔缩放和饱和程度，其作用需结合当前间隔与学习率分析。

DPO 优化两条回答相对参考策略的概率比。例如参考在三回答 $(w,l,o)$ 上为 $(0.4,0.4,0.2)$，某策略变为 $(0.3,0.1,0.6)$。优选概率从0.4降到0.3，但 $\frac{w}{l}$ 比从1升到3，因而 $\Delta=\log3>0$。因此应同时观察两条回答的概率、偏好间隔和自由生成表现。

取 $\beta=0.5$，上述间隔给出 $u=0.549306$，偏好概率为 $\frac{\sqrt3}{1+\sqrt3}\approx0.6340$，损失约0.4557，低于初始相等策略时的 $\log2\approx0.6931$。第三个回答在自由生成中的占比增大，其质量需要另行评价。表\ref{tab:rev-alignment-losses}总结了以 DPO 为代表的主流直接偏好优化损失及其工程技术特征。

\begin{table}[htbp]
\centering
\small
\setlength{\tabcolsep}{3.5pt}
\caption{主流偏好对齐与直接优化算法（Direct Alignment Losses）对比}
\label{tab:rev-alignment-losses}
\begin{tabularx}{\linewidth}{lp{18mm}cp{26mm}X}
\toprule
方法 & 数据格式 & 参考模型/显存 & 核心优化目标 & 核心优势与工程权衡 \\
\midrule
DPO & 成对 $(y_w, y_l)$ & 需参考（双模型） & $-\log \sigma(\beta \Delta_{\vect{\theta}})$ & 无需显式奖励模型与采样，训练稳定易收敛 \\
IPO & 成对 $(y_w, y_l)$ & 需参考（双模型） & $(\Delta_{\vect{\theta}} - \frac{1}{2\tau})^2$ & 施加平方正则项，防止偏好间隔发散过拟合 \\
KTO & 单样本 $(y, \pm)$ & 需参考（双模型） & 前景理论效用函数 & 仅需二元非配对标签，极大降低构建门槛 \\
ORPO & 成对 $(y_w, y_l)$ & 无需（单模型） & $\mathcal{L}_{\mathrm{SFT}} + \lambda \mathcal{L}_{\mathrm{odds}}$ & SFT 与偏好单阶段统一训练，彻底省去参考模型 \\
SimPO & 成对 $(y_w, y_l)$ & 无需（单模型） & 长度归一化奖励差加 $\gamma$ & 零参考模型显式长度惩罚，提升生成质量 \\
\bottomrule
\end{tabularx}
\end{table}

\subsection{序列概率与长度}
回答序列的条件概率必须遵循相同的停止协议：
\begin{equation}
 \log\pi_{\vect{\theta}}(y\mid x)=\sum_{t=1}^{T_y}\log\pi_{\vect{\theta}}(y_t\mid x,y_{<t}).
 \label{eq:pref-sequence-logprob}
\end{equation}
将求和除以 $T_y$ 得到平均词元对数概率，改变了回答之间的权重。长度归一化由此给出与式\eqref{eq:pref-regularized}不同的目标。

长度偏差可能同时来自标注偏好、候选生成和似然累积。参考比值抵消部分基线倾向，其余偏差仍在。若截断只保留两回答共同前缀，原本可辨别的比较可能变成无信号；若仅一条被截断，标签可能转而反映截断质量。共同提示应保持一致，优选与拒选的完整目标跨度应分别校验。

\begin{figure}[htbp]
\centering
\begin{tikzpicture}[node distance=8mm,box/.style={draw,rounded corners,align=center,text width=32mm,minimum height=11mm},>=stealth]
\node[box] (pair) {同一提示\\优选与拒选回答};
\node[box,right=of pair,yshift=10mm] (pol) {当前策略\\两条序列对数概率};
\node[box,right=of pair,yshift=-10mm] (ref) {冻结参考\\两条序列对数概率};
\node[box,right=of pol,yshift=-10mm] (diff) {相对间隔\\$\Delta_{\vect{\theta}}$};
\node[box,below=of diff] (loss) {$\operatorname{softplus}(-\beta\Delta_{\vect{\theta}})$\\仅更新策略};
\draw[->] (pair)--(pol);\draw[->] (pair)--(ref);\draw[->] (pol)--(diff);\draw[->] (ref)--(diff);\draw[->] (diff)--(loss);
\end{tikzpicture}
\caption{DPO 的数据与梯度边界。参考分支提供固定概率基线，全程不参与更新。}
\label{fig:pref-dpo}
\end{figure}

\section{在线 RLHF 的模型分工}

\term{参考策略}{Reference Policy}{}规定相对偏离的基准；\term{旧策略}{Behavior Policy}{}在本节指产生当前批次轨迹的行为策略。前者可以跨多次更新保持固定，后者随采样轮次变化。策略概率比用于控制复用样本时的更新，参考 KL 用于约束行为漂移，两者职责不同。

典型 Actor--Critic 链路中，策略生成回答，奖励模型给出序列评分，价值模型估计前缀状态的未来回报，参考模型提供 KL 基线。价值估计依赖当前策略和剩余过程，用于构造优势及降低估计方差；奖励评分表达对完整回答的评价。奖励模型可以在策略优化阶段冻结，但其输入分布会随策略变化。

序列 KL 可以用链式法则写成策略访问前缀下的条件 KL 之和：
\begin{equation}
 D_{\mathrm{KL}}(\pi_{\vect{\theta}}(y\mid x)\|\pi_{\mathrm{ref}}(y\mid x))
 =\mathbb E_{y\sim\pi_{\vect{\theta}}}\sum_t
 \log\frac{\pi_{\vect{\theta}}(y_t\mid x,y_{<t})}{\pi_{\mathrm{ref}}(y_t\mid x,y_{<t})}.
 \label{eq:pref-kl-chain}
\end{equation}
单个采样词元的对数比可以为负；在相应分布下取期望，才得到非负的 KL。若轨迹来自旧策略，直接平均这些值估计的是旧分布下的对数比，在作出近似或正确分布校正之前，该值只是当前策略 KL 的一个估计。

将终局奖励与逐词元 KL 惩罚结合，可构造供\bookxref{ch:reinforcement-learning}策略优化算法使用的回报。对数比中的当前策略参数也影响采样分布；实际代理目标通常在旧轨迹上冻结部分统计量并限制更新幅度，应把它当作近似优化过程，采样损失与式\eqref{eq:pref-regularized}的逐项精确梯度之间还有距离。

\begin{figure}[htbp]
\centering
\begin{tikzpicture}[node distance=8mm,box/.style={draw,rounded corners,align=center,text width=29mm,minimum height=10mm},>=stealth]
\node[box] (prompt) {版本化提示};
\node[box,right=of prompt] (policy) {策略生成\\保存行为概率};
\node[box,right=of policy] (reward) {奖励模型\\序列评分};
\node[box,below=of reward] (returns) {回报与优势\\KL、价值估计};
\node[box,left=of returns] (update) {受约束更新\\策略与价值模型};
\node[box,left=of update] (ref) {固定参考\\概率基线};
\draw[->] (prompt)--(policy);\draw[->] (policy)--(reward);\draw[->] (reward)--(returns);
\draw[->] (returns)--(update);\draw[->] (update)--(policy);\draw[->] (ref)--(returns);
\end{tikzpicture}
\caption{在线奖励优化的职责。参考、奖励、价值与策略是逻辑角色，不要求部署为四个完全独立的服务。}
\label{fig:pref-rlhf}
\end{figure}

\begin{algorithm}[离线 DPO 的一次参数更新]
\AlgoInput{一批偏好对、固定参考模型、策略、$\beta>0$，以及带版本的模板和分词器。}
\AlgoOutput{更新后的策略和目标统计量。}
\begin{myalgorithmic}[1]
\State Verify that each pair shares the same complete prompt, labels have defined semantics, and both response spans and termination protocols are valid; discard or separately record samples whose truncation destroys the preference basis
\State Compute target-token log-probability sums for both responses under the frozen reference; if cached, bind the key to reference weights, template, tokenizer, and exact token sequences
\State Compute log-probability sums over the same target positions under the current policy, then form the margin and numerically stable Softplus loss using Equation~\eqref{eq:pref-dpo}
\State Aggregate using the declared prompt or pair weights and backpropagate only into the policy; preserve the same weighting under accumulation and parallel reduction, noting that token averaging defines a different objective
\State Record loss, margin, both response log probabilities, lengths, and valid counts; periodically generate freely on independent prompts and compare blindly with the frozen baseline
\end{myalgorithmic}
\end{algorithm}

参考概率缓存只有在输入与参考完全固定时才有效。使用共享基座与适配器实现参考时，关闭某个适配器得到的未必是预期参考版本；还需确认基座未被更新、适配器组合一致、随机层处于确定的评分模式。

\begin{algorithm}[奖励模型与在线策略的分阶段更新]
\AlgoInput{比较数据、提示分布、初始策略及独立评估集。}
\AlgoOutput{带来源与评估记录的策略版本。}
\begin{enumerate}
\item 按提示和来源划分比较数据，拟合奖励模型，并评估排序、校准及关键切片的误差。
\item 固定本轮奖励和参考版本，从已标识的行为策略采样，保存回答、行为概率、终止原因和各奖励分量。
\item 依据声明的回报、优势和约束算法更新策略；限制样本复用范围，记录概率比、KL估计口径和裁剪比例。
\item 在独立标准下检查质量、长度、拒答和原有能力。若奖励增加而独立质量下降，停止扩大优化强度并分析候选分布漂移。
\item 若重新收集偏好或更新奖励模型，建立新版本与新的训练阶段；保留旧版本以便解释行为变化与恢复。
\end{enumerate}
\end{algorithm}

\section{GRPO 与方法选择}
\optional
\term{组相对策略优化}{Group Relative Policy Optimization}{GRPO}为同一提示生成 $G$ 条回答，以组内奖励构造相对优势\autocite{shao2024deepseekmath}。选择总体标准差约定时，
\begin{equation}
 \bar r=\frac1G\sum_i r_i,\qquad
 s=\sqrt{\frac1G\sum_i(r_i-\bar r)^2},\qquad
 A_i=\frac{r_i-\bar r}{s+\epsilon}.
 \label{eq:pref-grpo-overview}
\end{equation}
各优势之和为零；奖励全部相同时，相对优势全部为零。组内标准化消除绝对偏移，并重加权不同问题；标准化优势描述各回答在本组内的相对位置。

例如奖励 $(0,1,1,2)$ 的均值为1，标准差为 $\sqrt{\frac{1}{2}}$；忽略数值稳定项时优势为 $(-\sqrt2,0,0,\sqrt2)$。换成 $(10,10,10,10)$，尽管绝对奖励更高，组内没有排序信号。若奖励函数只识别格式，组内训练也只会提供格式差异。组内均值包含自身样本所带来的基线偏差、零方差组、长度归约和裁剪目标，在\bookxref{ch:verifiable-reasoning}中完整推导。

\begin{table}[htbp]
\centering
\caption{反馈学习方法的监督对象、计算角色与适用边界。}
\label{tab:preference-alignment-methods}
\begin{tabularx}{\linewidth}{@{}p{.13\linewidth}p{.24\linewidth}p{.27\linewidth}X@{}}
\toprule
方法 & 反馈对象 & 主要计算角色 & 关键边界\\
\midrule
SFT & 高质量示范回答 & 策略前后向 & 模仿示范；数据覆盖限制可学行为。\\
DPO & 离线成对偏好 & 策略与固定参考评分 & 无需每次更新在线生成，但受候选覆盖和偏好模型限制。\\
奖励模型加策略优化 & 比较数据及在线奖励 & 策略、奖励、参考；可有价值模型 & 采样与评分代价较高，奖励外推可能被利用。\\
GRPO & 同提示成组奖励 & 策略、评分器、可选参考约束 & 省去独立价值模型，但增加成组生成且可能无相对信号。\\
\bottomrule
\end{tabularx}
\end{table}

表\ref{tab:preference-alignment-methods}中的方法根据起始模型与反馈条件选择。是否需要示范微调取决于起始模型是否已有任务行为；是否需要在线采样取决于离线候选覆盖和评分可靠性。参数高效微调改变更新子空间，可与上述多种目标组合，反馈仍由示范、比较或奖励提供。

\section{奖励偏差与评估边界}

\term{奖励投机}{Reward Hacking}{}指策略利用评分机制的缺陷提高分数，却没有改善真正的任务目标。若代理奖励为 $\widehat r=r+e$，其中 $r$ 是目标效用、$e$ 是评分误差，则优化 $\widehat r$ 可以同时增加 $e$。原分布下误差均值很小，经过策略选择后的误差则可能显著增大。

偏好评估应同时固定候选生成预算、比较基线和呈现方式。胜率需要报告比较数、平局处理、提示分组及不确定性；重复比较同一提示并不等于独立扩大样本。长度、语言、任务类别和高严重度失败应分开观察，分别报告各类任务的表现与严重失败。

\begin{note}[对齐结论的适用范围]
奖励模型准确率、DPO间隔、在线奖励和独立任务质量分别测量不同对象。行为改进的评价应明确提示分布、评价标准、生成条件和参考基线。训练损失和奖励用于观察优化进展，独立任务与安全评估用于检验部署表现。
\end{note}




\chapterreferences
\myendchapterdocument
